\subsection{Common Unit}\label{subsec:common-unit}
The isotony axiom also contains this rather ``odd'' addendum:

\begin{quote}
We assume in addition that one of the two following situations prevails. Either $\mathfrak{U}(\mathbf{B_1})$ and $\mathfrak{U}(\mathbf{B_2})$ have a common unit element, or neither of them has a unit. The first situation can be obtained from the second by formal adjunction of a unit.
\end{quote}

If, as suggested, it is always possible, through formal adjunction, to add a unit, then why even consider the case in which there is no unit?

For example, using the state $\omega$ to apply the GNS construction to an (abstract) C*-algebra $\mathfrak{U}(\mathbf{B})$ that doesn't have a unit is impossible\footnote{Again, technically, one only requires the existence of a ``bounded, approximate identity'' and not a unit to produce a distinguished vector $\Omega_\omega$. However, physical motivation for the existence of a ``bounded, approximate identity'' as opposed to a unit is lacking.}. The presence of a unit in $\mathfrak{U}(\mathbf{B})$ leads directly to the distinguished vector $\Omega_\omega$ in the Hilbert space $\mathcal{H}_\omega$. This distinguished vector $\Omega_\omega$ is interpreted as being a vacuum. So having a unit in $\mathfrak{U}(\mathbf{B})$ seems indispensable.

That being said, we will simply assume that all such $\mathfrak{U}(\mathbf{B})$ contain a unit.

So if we adopt the position that a unit is always present, one convenient way to include it, and the method we adopt, is to assume that the empty set $\emptyset$ is assigned the C*-algebra $\mathbb{C} \mathbf{1}$ generated by the unit $\mathbf{1}$
\begin{align}
\emptyset \rightarrow \mathfrak{U}(\emptyset) \equiv \mathbb{C} \mathbf{1}.
\end{align}
As the empty set $\emptyset$ is a member of any $\mathbf{B}$, isotony implies $\mathfrak{U}(\emptyset) \subset \mathfrak{U}(\mathbf{B})$ which in turn implies that all $\mathfrak{U}(\mathbf{B})$ contain a unit.
